\documentclass[10pt,a4paper]{amsart}
\usepackage[margin=19mm]{geometry}
\usepackage{amsmath,amssymb,mathtools}
\usepackage[scr=rsfs,cal=euler]{mathalfa}
\usepackage{xfrac}
\usepackage[unicode,hidelinks,bookmarks=false]{hyperref}
\newcommand{\ie}{i.e.\ }
\newcommand{\cf}{cf.\ }
\newcommand{\resp}{respectively\ }
\newcommand{\Subsection}{\subsection}
\providecommand{\nobreakdash}{\nobreak\hbox{-}}
\setlength{\emergencystretch}{2em}
\multlinegap0pt
\newtheorem{lem}{Lemma}
\newtheorem{thm}{Theorem}
\title{Normal Equations and Discrete Energy Structures for High-Order Streamline Diffusion}
\author{Erik Burman, Peter Hansbo, Mats G. Larson}
\date{Source: arXiv:2609.12506v1 (CC BY 4.0)}
\begin{document}
\maketitle

\noindent\emph{Source and licence.} Normal Equations and Discrete Energy Structures for High-Order Streamline Diffusion. Version arXiv:2609.12506v1 (CC BY 4.0), distributed by arXiv under the Creative Commons Attribution 4.0 International licence, http://creativecommons.org/licenses/by/4.0/, as recorded in the arXiv licence metadata for this exact version and shown on the abstract page https://arxiv.org/abs/2609.12506v1. Full licence notice: \texttt{licenses/arxiv-2609.12506-CC-BY-4.0.txt}.\par\smallskip

\noindent\textit{Editorial scope.} Counted unit: the article's thm thm:adams-moulton-stability (source0.tex lines 1742--1763). The article's complete original proof of it is delivered with it (source0.tex lines 1764--1827, one \texttt{proof} environment of 64 lines). No mathematics is written, repaired or re-derived: the statement and the proof are the article's own lines, in the article's own notation, and no retained line is substituted or re-flowed.  Retained uncounted as the source-local context the proof needs: lines 1229--1232; lines 1234--1237; lines 1242--1245; lines 1268--1273; lines 1530--1582; lines 1625--1645; lines 1655--1672; lines 1659--1662.  Nothing was omitted from any retained span, so the package carries no omission record.  No retained line contains a citation, so the wrapper carries no bibliography and no bibliography entry.  No retained line contains a figure environment, an image-inclusion command or drawing code, so no image is delivered and the package declares no asset dependency.  Editorial formatting only, in addition: an AMS \texttt{amsart} preamble that restates the article's own declarations and macros used by the retained lines, namely the environments \texttt{lem}, \texttt{thm} on separate counters, together with the standard packages those lines need.  The retained environments print in this document as Lemma 1, Theorem 1 in order of appearance, whereas the article prints the same items with the numbering of its own full document.  The complete native source of the article as distributed in the arXiv e-print for this exact version is bundled unchanged as \texttt{sources/210121/source0.tex}.
\begin{equation}
\|GV_h\|_{\Omega}\le C_Gh^{-1}\|V_h\|_{\Omega}
\quad \forall V_h\in\mathcal V_h \label{eq:adams-inverse-inequality}
\end{equation}

\begin{equation}
0<\delta\le C_{\delta}h,\qquad
\tau\le c_0C_G^{-1}h \label{eq:cflcond}
\end{equation}

\begin{align}
\tau&\le C_{4/3}h^{4/3} &&\text{for AM3} \label{eq:am3-strengthened-cfl}\\
\tau&\le C_{6/5}h^{6/5} &&\text{for AM4} \label{eq:am4-strengthened-cfl}
\end{align}

\begin{equation}
(A_{\tau}V_h^n,T_{\delta}W)_{\Omega}
=(R^n,T_{\delta}W)_{\Omega}
\quad\forall W\in\mathcal V_h,\qquad n\ge n_0(\nu)
\label{eq:am-general-forced-equation}
\end{equation}

\begin{lem}[Exact Adams--Moulton decompositions]
\label{lem:adams-exact-energy-decompositions}
For \(\nu=3,4,5\),
\begin{equation}
\mathcal E_N^{(\nu)}=\mathcal B_N^{(\nu)}+\mathrm{BT}_N^{(\nu)} \label{eq:adams-bulk-boundary-decomposition}
\end{equation}
where the bulk and boundary terms are as follows.

\paragraph{AM3.}
\begin{align}
\mathcal B_N^{(3)}
&=-\frac1{24}\sum_{n=2}^N\|d^2U^n\|_{\delta}^2 \label{eq:am3-energy-bulk}\\
\mathrm{BT}_N^{(3)}
&=\frac12\bigl(\|U^N\|_{\delta}^2-\|U^1\|_{\delta}^2\bigr)
-\frac1{24}\bigl(\|d^1U^N\|_{\delta}^2
-\|d^1U^1\|_{\delta}^2\bigr) \label{eq:am3-energy-boundary}
\end{align}

\paragraph{AM4.}
\begin{align}
\mathcal B_N^{(4)}
&=-\frac1{48}\sum_{n=3}^N\|d^3U^n\|_{\delta}^2 \label{eq:am4-energy-bulk}\\
\mathrm{BT}_N^{(4)}
&=\frac12\bigl(\|U^N\|_{\delta}^2-\|U^2\|_{\delta}^2\bigr)
-\frac1{24}\bigl(\|d^1U^N\|_{\delta}^2
-\|d^1U^2\|_{\delta}^2\bigr) \label{eq:am4-energy-boundary-main}\\
&\quad-\frac1{24}\bigl((d^1U^N,d^2U^N)_{\delta}
-(d^1U^2,d^2U^2)_{\delta}\bigr)
+\frac1{48}\bigl(\|d^2U^2\|_{\delta}^2
-\|d^2U^N\|_{\delta}^2\bigr) \label{eq:am4-energy-boundary-increments}
\end{align}

\paragraph{AM5.}
\begin{align}
\mathcal B_N^{(5)}
&=\frac3{160}\sum_{n=4}^{N-1}\|d^3U^n\|_{\delta}^2
-\frac{19}{1440}\sum_{n=4}^{N}\|d^4U^n\|_{\delta}^2 \label{eq:am5-energy-bulk}\\
\mathrm{BT}_N^{(5)}
&=\frac12\bigl(\|U^N\|_{\delta}^2-\|U^3\|_{\delta}^2\bigr)
-\frac1{24}\bigl(\|d^1U^N\|_{\delta}^2
-\|d^1U^3\|_{\delta}^2\bigr) \label{eq:am5-energy-boundary-main}\\
&\quad-\frac1{24}(d^1U^N,d^2U^N)_{\delta}
+\frac1{24}(d^1U^3,d^2U^3)_{\delta} \notag\\
&\quad-\frac1{48}\|d^2U^N\|_{\delta}^2
+\frac1{48}\|d^2U^3\|_{\delta}^2
+\frac{19}{1440}\|d^2U^{N-1}\|_{\delta}^2
-\frac{19}{1440}\|d^2U^2\|_{\delta}^2 \notag\\
&\quad-\frac{19}{720}(d^1U^N,d^3U^N)_{\delta}
+\frac{19}{720}(d^1U^3,d^3U^3)_{\delta} \notag\\
&\quad-\frac{11}{1440}\|d^3U^N\|_{\delta}^2
+\frac{19}{720}\|d^3U^3\|_{\delta}^2 \label{eq:am5-energy-boundary-increments}
\end{align}
\end{lem}

\begin{lem}[Bulk lower bounds]
\label{lem:adams-bulk-lower-bounds}
Let \(V_h\) solve \eqref{eq:am-general-forced-equation}, and evaluate the
bulk terms on this sequence. Under \eqref{eq:cflcond} and the relevant
strengthened condition, the AM3 and AM4 bulk terms satisfy
\begin{align}
\mathcal B_m^{(\nu)}
&\ge-C\tau C_{2(\nu-1)/(2\nu-3)}^{\,2\nu-3}
\sum_{n=0}^{m-1}\|V_h^n\|_{\delta}^2 \label{eq:am34-bulk-lower-bound-solution}\\
&\quad-C\tau^2\sum_{n=n_0(\nu)}^{m}\|R^n\|_{\Omega}^2
-C\mathcal I_{\tau}(V_h) \label{eq:am34-bulk-lower-bound-forcing}
\end{align}
For AM5,
\begin{align}
\mathcal B_m^{(5)}
&\ge\left(\frac3{160}-Cc_0^2\right)
\sum_{n=4}^{m-1}\|d^3V_h^n\|_{\delta}^2 \label{eq:am5-bulk-lower-bound-solution}\\
&\quad-C\tau^2\sum_{n=n_0(5)}^{m}\|R^n\|_{\Omega}^2
-C\mathcal I_{\tau}(V_h) \label{eq:am5-bulk-lower-bound-forcing}
\end{align}
\end{lem}

\begin{lem}[Boundary lower bound]
\label{lem:adams-boundary-lower-bound}
Let \(V_h\) solve \eqref{eq:am-general-forced-equation}. For every truncated
horizon \(J\), \(n_0(\nu)\le J\le N\), choose \(m_J\) so that
\begin{equation}
\|V_h^{m_J}\|_{\delta}
=\max_{n_0(\nu)\le n\le J}\|V_h^n\|_{\delta} \label{eq:adams-maximum-energy-index}
\end{equation}
For \(c_0\) sufficiently small,
\begin{align}
\mathrm{BT}_{m_J}^{(\nu)}
&\ge\left(\frac12-Cc_0^2\right)
\|V_h^{m_J}\|_{\delta}^2
-C\mathcal I_{\tau}(V_h) \label{eq:adams-boundary-lower-bound-solution}\\
&\quad-C\tau^2\sum_{n=n_0(\nu)}^{m_J}\|R^n\|_{\Omega}^2
\label{eq:adams-boundary-lower-bound-forcing}
\end{align}
\end{lem}

\begin{equation}
\|V_h^{m_J}\|_{\delta}
=\max_{n_0(\nu)\le n\le J}\|V_h^n\|_{\delta} 
\end{equation}

\begin{thm}[Adams--Moulton stability]
\label{thm:adams-moulton-stability}
Assume the inverse inequality \eqref{eq:adams-inverse-inequality} and the CFL
condition \eqref{eq:cflcond}, with \(c_0\) sufficiently small. For AM3 also
assume \eqref{eq:am3-strengthened-cfl}, and for AM4 assume
\eqref{eq:am4-strengthened-cfl}. If \(V_h\) solves
\eqref{eq:am-general-forced-equation}, then
\begin{equation}
\|V_h\|_{\mathcal E_\nu,N}^2
\lesssim\mathcal I_{\tau}(V_h)
+\|R\|_{\mathcal E_\nu^*,N}^2
+\tau\|R\|_{\ell_{\tau}^2(n_0(\nu),N;\mathcal H)}^2
\label{eq:adams-stability-right-hand-side}
\end{equation}
For AM3 and AM4, the hidden constant depends exponentially on
\(TC_{2(\nu-1)/(2\nu-3)}^{\,2\nu-3}\). For AM5 it is independent of a
strengthened CFL constant.
The theorem requires only \(0<\delta\le C_{\delta}h\); hence all three
conclusions include the symmetric positive definite normal-equation choice
\(\delta=b_0\tau\) under
their stated CFL conditions.
\end{thm}

\begin{proof}
The proof has two distinct branches. For AM3 and AM4, the negative bulk term
is controlled by the increment recurrence and then by a discrete Gronwall
argument. For AM5, the positive third-increment term absorbs the negative
fourth-increment term, so no Gronwall step is needed.
Use \(S_{\delta}V_h^n\) as the test function in
\eqref{eq:am-general-forced-equation}, sum to an arbitrary horizon, and
insert Lemma \ref{lem:adams-exact-energy-decompositions}. Write
\begin{equation}
K_\nu=C_{2(\nu-1)/(2\nu-3)}^{\,2\nu-3},\qquad
M_J=\max_{n_0(\nu)\le n\le J}\|V_h^n\|_{\delta}^2
\label{eq:am-gronwall-quantities}
\end{equation}
For each \(J\), choose the running maximum index \(m_J\) from
\eqref{eq:adams-maximum-energy-index} and apply the energy identity on the
truncated interval ending at \(m_J\). For AM3 and AM4, Lemmas
\ref{lem:adams-bulk-lower-bounds} and
\ref{lem:adams-boundary-lower-bound}, together with the dual-norm definition
and Young's inequality, give
\begin{align}
M_J
&\le C\mathcal A_N
+C\tau K_\nu\sum_{n=n_0(\nu)}^{J-1}M_n
\label{eq:am-running-maximum-recursion}\\
\mathcal A_N
&=\mathcal I_{\tau}(V_h)+\|R\|_{\mathcal E_\nu^*,N}^2
+\tau\|R\|_{\ell_{\tau}^2(n_0(\nu),N;\mathcal H)}^2
\label{eq:am-stability-data-quantity}
\end{align}
Crucially, the bulk estimate contains
\(\sum_{n=0}^{m_J-1}\|V_h^n\|_{\delta}^2\), not the current value. Its
start-up part is controlled by \(\mathcal I_{\tau}(V_h)\), and its computed
part is bounded by \(\sum_{n=n_0(\nu)}^{J-1}M_n\). The dual forcing term is
bounded by
\begin{equation}
\|R\|_{\mathcal E_\nu^*,N}
\bigl(\mathcal I_{\tau}(V_h)
+\|V_h\|_{\mathcal E_\nu,m_J}^2\bigr)^{1/2}
\label{eq:am-dual-forcing-at-running-maximum}
\end{equation}
so Young's inequality absorbs the running graph energy and the partial
material-residual sum. The endpoint and bulk forcing defects are exactly
\(C\tau^2\sum_{n=n_0(\nu)}^{m_J}\|R^n\|_{\Omega}^2\), which is bounded by the
last term of \(\mathcal A_N\).

The standard discrete Gronwall lemma applied to
\eqref{eq:am-running-maximum-recursion} yields
\begin{equation}
M_N\le C\mathcal A_N\exp(C T K_\nu)
\label{eq:am-running-maximum-gronwall}
\end{equation}
which displays the claimed dependence on the strengthened CFL constant.
Applying the energy identity once more at the final horizon \(N\), using
\eqref{eq:am-running-maximum-gronwall} for all preceding solution terms, and
again applying Young's inequality controls the complete
\(\tau\delta\sum_{n=n_0(\nu)}^N\|A_{\tau}V_h^n\|_{\Omega}^2\) sum.

For AM5, choose \(c_0\) so that the negative fourth-increment contribution
in Lemma \ref{lem:adams-bulk-lower-bounds} is absorbed by the positive
third-increment sum. The same running-maximum and dual-norm argument then
controls the graph energy without Gronwall. Repeating it at \(N\) controls
the full material-residual and third-increment sums. These arguments prove
\eqref{eq:adams-stability-right-hand-side} for the general forcing \(R\).
\end{proof}

\end{document}
